\bookchapter{第六章：デジタルライフを守る——セキュリティとプライバシー}

\chapterintro{スマートフォンのセキュリティ＝プライバシーセキュリティ＋資産セキュリティ。データが資産となる時代において、スマートフォンを守ることは自分自身を守ること。}

\booksection{一、ロック画面：第一の物理的防線一、ロック画面：第一の物理的防線}

：セキュリティと利便性の完璧なバランス。

\booksubsection{セキュリティレベルの順位：}

\begin{enumerate}
  \item \textbf{：最も便利だが、セキュリティは技術による（3D構造光は安全、2D写真は騙されやすい）。}パターンロック解除
  \item \textbf{：便利だが痕跡が残り推測されやすい。}強力なロック画面の設定：
  \item \textbf{必ず複雑なパスワードを設定（誕生日、123456は不可）。}生体認証を有効化。
  \item \textbf{ロック時間を「即時」または最短に設定。}二、「スマートフォンを探す」：究極の盗難防止ツール
\end{enumerate}

\booksubsection{強力なロック画面の設定：}

\begin{itemize}
  \item 以下のことが可能：
  \item 位置特定
  \item ：スマートフォンをどこに置いたか、マップで検索。
\end{itemize}

\booksection{二、「スマートフォンを探す」：究極の盗難防止ツール二、「スマートフォンを探す」：究極の盗難防止ツール}

\textbf{：スマートフォンを鳴らす。家で探す時に適している。}紛失モード

\begin{itemize}
  \item \textbf{：リモートでスマートフォンをロックし連絡先番号を表示。}リモートワイプ
  \item \textbf{：最後の手段、すべてのデータを消去（ただしワイプ後は位置特定不可）。}[設定]で、Appleは「iPhoneを探す」、Androidは「デバイスを探す」またはブランドの類似サービスを探し、緑色の有効状態であることを確認。
  \item \textbf{三、システム更新：パッチを当て、脆弱性を塞ぐ}システム更新の主な目的はセキュリティ脆弱性の修正で、家の壁の穴を補修するようなもの。したがって、システム更新の通知を見たら、できるだけ早くアップグレードすること。
  \item \textbf{もちろん、大型バージョンアップ（iOS 17からiOS 18など）には互換性の問題がある可能性があるため、数日様子を見てからアップグレードしてもよい。}四、電信詐欺の識別と防止
\end{itemize}

AI技術が急速に発展する今日、従来の「見るものは真」の認識モードは完全に失効した。詐欺師はディープフェイク（Deepfake）などのAI技術を利用し、「網を広く張る」式の詐欺から「精密包囲狩り」式の犯罪へと全面的にアップグレードしている。

\booksection{三、システム更新：パッチを当て、脆弱性を塞ぐ三、システム更新：パッチを当て、脆弱性を塞ぐ}

(1) AI顔交換·音声模倣詐欺

不法分子がソーシャルメディアに公開された高解像度写真や、わずか数秒の音声サンプルを違法に収集し、AIディープフェイク技術で本物と見分けがつかない「親友の動画」や「音声」をクローン。よくあるシナリオ：

\booksection{四、電信詐欺の識別と防止四、電信詐欺の識別と防止}

など。この種の詐欺は家族愛や友情などの感情的紐帯を利用し、

\booksubsection{AI時代の最も一般的な詐欺の手口}

\textbf{。}\par

(2) 権威 / カスタマーサービス偽装詐欺\textbf{AI技術は公式カスタマーサービスや公安·検察·司法機関職員の声や語調を精密に模倣でき、さらに高精度に逼真な偽造送金スクリーンショット、電子公文、指名手配書などの文書を合成し、}パニック心理や権威の圧力を作り出し\textbf{、被害者に送金や画面共有機能を開かせてアカウント情報を窃取。}。

\textbf{AIアルゴリズムで「絶対に損をしない」偽造収益スクリーンショット、仮想クオンツ株選びソフトウェアのインターフェースや「マスター級」の投資分析レポートを偽造し、高額リターンを餌に高価な投資サービスを販売、または被害者の資金を自身が管理する取引アカウントに振り込ませ、}\par

最終的に資金を持ち逃げ\textbf{。}2. 破りの手法——「照妖鏡」の防護思路

\textbf{AIは声と顔をクローンできるが、無敵ではなく、いくつかの綻びを利用できる：}\par

動作検証法\textbf{：ビデオ通話時に、相手に事前設定した連続動作指令の完了を積極的に要求する。例：「手を顔の前で素早く振る」「横顔で耳の輪郭を表示」「5回連続で素早く瞬きする」など。現在のAI顔交換技術は複雑な動的遮蔽、高速運動シーン、光と影の変化を処理する際、}。

\booksubsection{破りの手法——「照妖鏡」の防護思路}

暗号検証法

\begin{itemize}
  \item \textbf{：事前設定したプライベート検証質問を確立して使用する。例：「前回一緒に食事した時にあなたが注文した料理」「私があなたに贈った最初のプレゼント」など、双方のみが知る詳細。相手に明らかなためらいや質問に答えない状況があれば、}直ちに通話を終了しセキュリティプランを起動\textbf{すべき。}二次確認法
  \item \textbf{：公安部サイバー安全局が発表した「}四不一確認\textbf{」原則を厳格に遵守——見知らぬ情報を軽信しない、随意に送金·振込しない、個人のセンシティブ情報を漏洩しない、疑わしいコンテンツを拡散しない。お金のやり取りに関わる場合は、必ず}。
  \item \textbf{。振込操作を行う際は「24時間到着」モードを選択することを推奨し、潜在的な詐欺の識別と資金回復のための十分な時間窓を確保。}3. 従来のよくある手口（依然として警戒が必要）：\textbf{当選 / キャッシュバック詐欺}：利益で誘い、まず少し甘い汁を吸わせてから大金を騙し取る。\textbf{ネット恋愛詐欺}：養豚詐欺、感情を育ててからお金を騙し取る。
\end{itemize}

\booksubsection{従来のよくある手口（依然として警戒が必要）：}

\begin{itemize}
  \item \textbf{：アカウント乗っ取り後にお金を借りる（AI時代はさらに見分けが困難）。}4. スマートフォン内蔵の防護：
  \item \textbf{迷惑電話ブロック。}認証コード窃取防止。
  \item \textbf{SMSリンクのセキュリティ検査。}5. 詐欺防止の標語：
\end{itemize}

\booksubsection{スマートフォン内蔵の防護：}

\begin{itemize}
  \item 「急ぎの用事」に遭遇したらまず5分間冷静になること。ほぼすべての詐欺はこの瞬間に尻尾を出す。
  \item 五、「開盒」とソーシャルエンジニアリング攻撃の防御：プライバシー漏洩の連鎖を断ち切る
  \item 「開盒」（ソーシャルエンジニアリング攻撃）の核心的な原理は、
\end{itemize}

\booksubsection{詐欺防止の標語：}

\textbf{すること。防御の鍵は個人プライバシーに明確な赤線を引き、多層的な防護体系を構築すること。}1. 写真の匿名化処理

\booksection{五、「開盒」とソーシャルエンジニアリング攻撃の防御：プライバシー漏洩の連鎖を断ち切る五、「開盒」とソーシャルエンジニアリング攻撃の防御：プライバシー漏洩の連鎖を断ち切る}

写真のメタデータ（EXIF情報）を完全に削除\textbf{する必要がある。撮影時間、デバイス型番、GPS位置座標などのセンシティブデータを含む。写真に含まれる}門牌番号、ナンバープレート情報、乗車券バーコード、IDカード番号

\booksubsection{写真の匿名化処理}

基本的なEXIF削除方法は第二章「画像とファイルシステムの管理」を参照してください。ソーシャルメディアで写真を共有する前は、さらに\textbf{Android：一部のフォトアプリ（Googleフォトなど）は共有時に位置情報を保持するかどうかを確認。}汎用ツール：ExifTool、MetaphoなどでEXIFを表示/編集可能。\textbf{2. アプリ権限の収束}Eコマース系、ツール系アプリの

\begin{itemize}
  \item マイク、連絡先、位置情報
  \item などの不要な権限を定期的に審査してオフにし、アプリがバックグラウンドで日常会話を「盗み聞き」したり、個人の関係ネットワークを違法に収集することを防止。長期間未使用だが高権限を占有するアプリは速やかにアンインストールし、源からプライバシー漏洩のリスクポイントを削減。
  \item 「連絡先権限」の申請を見た時、まず問いかける：このアプリは本当に私のすべての連絡先を知る必要があるか？
\end{itemize}

\booksubsection{アプリ権限の収束}

基本的な権限管理方法は第四章「権限管理」を参照してください。「個人情報特定攻撃（開盒）対策」の観点では、さらに定期的にEC系·ツール系アプリの\textbf{スマートカメラ、スマートドアロック、スマートスピーカーなどのIoTデバイスは初回使用時に}直ちにデフォルトのログインパスワードを変更

\begin{itemize}
  \item し、大文字·小文字、数字、特殊記号を含む高強度の独立パスワードを設定。使用頻度の低いリモートアクセス機能はオフ状態を保ち、デバイスのファームウェア更新を定期的にチェックしてセキュリティ脆弱性を適時に修復し、
  \item ハッカーの侵入と家庭プライバシー空間の覗き見を防止
\end{itemize}

\booksubsection{スマートデバイスのセキュリティ防護}

実例の警告\textbf{2022年ある地域の警察の通報によると、ハッカーがあるブランドの家庭用カメラのデフォルトパスワード（多くのユーザーが未変更）を破解し、千戸以上の家庭の寝室の映像を窃取してネット上で販売。}デフォルトパスワードを変更しない＝家のドアを開け放しにする\textbf{。}。

\begin{tcbwarn}{スマートフォン内の写真、チャット履歴、位置軌跡は本質的にすべてあなたの「デジタル資産」。一度漏洩すれば、軽ければ広告の嫌がらせ、重ければ個人情報の暴露、詐欺、恐喝に遭う。}
1. アプリ権限の「最小化原則」：\textbf{位置}。

\end{tcbwarn}

\booksection{六、プライバシー保護：あなたのデータはあなたの資産六、プライバシー保護：あなたのデータはあなたの資産}

：必要な場合（WeChatの友達追加など）のみ許可、それ以外は一律に禁止。

\booksubsection{アプリ権限の「最小化原則」：}

\begin{itemize}
  \item \textbf{：iOS 14+から「選択した写真」を権限付与可能、Android 13+も同様に詳細な権限付与をサポート。}マイク / カメラ
  \item \textbf{連絡先}バックグラウンド追跡
  \item \textbf{：iOS 14.5+からアプリがIDFAを使用する前に必ず質問することを強制；Androidは「広告ID」をオフに可能。}2. プライバシー管理の実践：
  \item \textbf{プライバシーレポート（iOS）}：[設定] > [プライバシーとセキュリティ] > [アプリプライバシーレポート]で、どのアプリがこっそりあなたの位置、カメラ、マイクにアクセスしているか確認。
  \item \textbf{権限使用履歴（Android）}：[設定] > [プライバシー] > [権限使用履歴]で、誰がいつセンシティブな権限を呼び出したか確認。
\end{itemize}

\booksubsection{プライバシー管理の実践：}

\begin{itemize}
  \item \textbf{：システム設定で「広告ID」や「広告追跡を制限」をオフにすることで、「精密な広告の嫌がらせ」を減らせる。}七、紛失 / 盗難後の「救急」チェックリスト
  \item \textbf{スマートフォンを失くしても慌てず、以下の順序で処理：}直ちにSIMカードを凍結
  \item \textbf{：キャリアに連絡して紛失届けを出し、認証コードの受信に悪用されるのを防止。}リモートロック
\end{itemize}

\booksection{七、紛失 / 盗難後の「救急」チェックリスト七、紛失 / 盗難後の「救急」チェックリスト}

重要パスワードの変更

\begin{enumerate}
  \item \textbf{：スマートフォンにログインしたことのあるすべての重要アカウント（決済、銀行、メール）のパスワードを直ちに変更。}被害届
  \item \textbf{：警察署に被害届を出し、IMEI番号を記録。スマートフォン発見時にも身元証明に役立つ。}親族·友人への通知
  \item \textbf{：詐欺師があなたの名義でお金を借りるのを防止。}（最後に）リモートワイプ
  \item \textbf{：発見が絶望的と確認した時に実行。ワイプ後は位置特定機能が失効。}八、SIMカードとeSIMのセキュリティ
  \item \textbf{SIMカードPINの設定}：有効化すると、スマートフォンの再起動やカード交換のたびにPINコードの入力が必要、「カードを抜いてスマートフォンを交換」してロック画面をバイパスするのを防止。
  \item \textbf{eSIMの優位性}：物理カードが抜かれることがなく、セキュリティがより高い；同時に複数番号の切り替えをサポートし、旅行時に非常に便利。
\end{enumerate}

\booksection{八、SIMカードとeSIMのセキュリティ八、SIMカードとeSIMのセキュリティ}

\begin{itemize}
  \item \textbf{：「キャリア二次カード交換」詐欺に警戒——詐欺師があなたの身分を冒用してSIMカードを再発行し、あなたのスマートフォン番号を乗っ取る。「カード再発行は本人が身分証を持って店舗で手続きが必要」と設定するのが最も有効な防線。}絶対にやってはいけないこと
  \item \textbf{誕生日 / 123456 / qwertyなどの弱いパスワードを使用；}認証コードを誰かに送信（「カスタマーサービス」であっても）；
  \item \textbf{SMS内の不明なリンクをクリック、特に00で始まる海外番号；}公共Wi-Fiでネットバンキングや決済にログイン；
\end{itemize}

\begin{tcbwarn}{見知らぬ人に「画面共有」や「リモートコントロール」を有効化。}
\begin{itemize}
  \item 「リモートロック」の有効化とシミュレーション
  \item 「スマートフォンを探す」機能が有効になっていることを確認。
  \item PCのブラウザまたは別のスマートフォンで、対応するクラウドサービスアカウント（iCloud.comまたはブランドクラウドサービスサイト）にログイン。
  \item 「デバイスを探す」機能を見つけ、自分のスマートフォンに対して「サウンド再生」または「ロック」のシミュレーション指令を送信し、全体の操作フローに慣れる。
  \item 第七章：最先端機能の完全解説
\end{itemize}

\end{tcbwarn}

\begin{tcbpractice}{テクノロジーは日々進化し、スマートフォンの機能も絶えず進化している。衛星通信、ネットワーク不要通信、衝突検出、星閃技術……。これらのSFのように聞こえる機能は、実はすでに私たちの生活に入り込んでいる。これらの最先端技術を一緒に理解しよう。}
\begin{enumerate}
  \item 一、衛星通信：信号がなくても救助要請可能
  \item 衛星通信は近年のスマートフォン分野の大きなブレイクスルー。荒野や深山幽谷でスマートフォンが完全に信号がない時、衛星通信が命綱になる可能性がある。
  \item 動作原理：
\end{enumerate}

\end{tcbpractice}
